%-------------------------
% Cover Letter in Latex
% Enrique Martínez Martel
%------------------------

\documentclass[letterpaper,11pt]{article}

\usepackage{latexsym}
\usepackage[empty]{fullpage}
\usepackage{marvosym}
\usepackage[usenames,dvipsnames]{color}
\usepackage[hidelinks]{hyperref}
\usepackage{fancyhdr}
\usepackage[english]{babel}
\usepackage{tabularx}
\usepackage{ragged2e}
\usepackage{fontawesome5}
\usepackage{graphicx}
\usepackage[scale=0.90,lf]{FiraMono}
\usepackage{enumitem}

\definecolor{light-grey}{gray}{0.83}
\definecolor{dark-grey}{gray}{0.3}
\definecolor{text-grey}{gray}{.08}

\usepackage{contour}
\usepackage[normalem]{ulem}
\renewcommand{\ULdepth}{1.8pt}
\contourlength{0.8pt}
\newcommand{\myuline}[1]{%
  \uline{\phantom{#1}}%
  \llap{\contour{white}{#1}}%
}

\usepackage{tgheros}
\renewcommand*\familydefault{\sfdefault}
\usepackage[T1]{fontenc}

\pagestyle{fancy}
\fancyhf{}
\fancyfoot{}
\renewcommand{\headrulewidth}{0pt}
\renewcommand{\footrulewidth}{0pt}

\addtolength{\oddsidemargin}{-0.5in}
\addtolength{\evensidemargin}{0in}
\addtolength{\textwidth}{1in}
\addtolength{\topmargin}{-.8in}
\addtolength{\textheight}{1.5in}

\urlstyle{same}
\raggedbottom
\raggedright

\color{text-grey}

\begin{document}

\vspace*{160pt}

Dear Prof. Vechev,

\vspace{16pt}

I am writing to express my interest in joining the SRI Lab as a PhD student. My background is in deep learning with a focus on transformer architectures and large language models: a multimodal LLM classification pipeline built at Banco Santander evaluating GPT-4/GPT-4o on real operational data, an MSc thesis at UPM benchmarking deep learning architectures on a real-world dataset (grade: 10/10), and experience developing and maintaining ML pipelines in production at scale. More of my work is at \href{https://enkairito.github.io/enriquemm.github.io/}{\myuline{enkairito.github.io/enriquemm.github.io}}.

\vspace{16pt}

The main reasons I want to do this PhD:

\vspace{8pt}

\begin{itemize}[leftmargin=1.5em, itemsep=10pt]

  \item \textbf{Reliable LLM evaluation is itself an open research problem.} My MSc thesis was a comparative analysis of computer vision techniques on a custom real-world dataset, designed to understand generalization, failure modes, and deployment trade-offs across architectures, not just to report a number on a standard split. That experience made clear how much the methodology shapes the conclusion. The Reliable Language Model Evaluation project asks exactly that question at scale, and it is where I think I can contribute most directly.

  \item \textbf{LLMs for code require constraints that general language modeling ignores.} At Banco Santander I built and evaluated a multimodal LLM pipeline using GPT-4/GPT-4o on real operational data, including benchmarking of prompting strategies and model behavior under distribution shift. Code generation is a harder version of the same problem: the output has to be not just plausible but correct, and increasingly, secure. Constrained decoding and adversarial robustness in that setting are problems I want to work on.

  \item \textbf{Research that produces things.} DeepCode, ChainSecurity, LatticeFlow: the SRI Lab has a track record of turning foundational ideas into systems that people actually use. That is the kind of research I want to do. Not papers that get cited, but work that changes how something is done.

\end{itemize}

\vspace{16pt}

I do not have a programming languages or formal methods background. I am not pretending otherwise. But I have built and deployed ML systems at scale, I understand transformers from the inside, and I learn new domains fast when the problem is worth understanding. I think my applied side complements the theoretical depth the lab already has.

\vspace{16pt}

I would be happy to discuss further.

\vspace{20pt}

Best regards,

\vspace{16pt}

\textbf{Enrique Martínez Martel}

\vspace{6pt}

\faEnvelope \hspace{2pt} \texttt{enri.martinez25@gmail.com} \hspace{6pt} $|$ \hspace{6pt}
\faLinkedin \hspace{2pt} \href{https://www.linkedin.com/in/enrique-martinez-martel}{\myuline{enrique-martinez-martel}} \hspace{6pt} $|$ \hspace{6pt}
\faGlobe \hspace{2pt} \href{https://enkairito.github.io/enriquemm.github.io/}{\myuline{enkairito.github.io/enriquemm.github.io}}

\end{document}
